\part{Construcción, evaluación y mejora de soluciones}

\chapter{Construcción de soluciones mediante fundamentos de programación}
\label{cap:03-fundamentos-de-programacion}
\chaptermark{Fundamentos de programación}

Los capítulos anteriores establecieron las bases para comprender un problema y diseñar su solución. En este capítulo se transforma ese diseño en instrucciones ejecutables. Se utiliza Python 3 como lenguaje de aprendizaje, manteniendo el vínculo entre el modelo de entrada, procesamiento y salida, el algoritmo y su implementación. El objetivo no es memorizar instrucciones aisladas, sino comprender cómo se combinan para construir una solución coherente.

\begin{resultados}
  \begin{itemize}
    \item Traducir algoritmos a código Python y seleccionar tipos de datos y operadores.
    \item Implementar secuencias, decisiones y ciclos.
    \item Definir funciones, manipular colecciones y comprobar resultados con datos de prueba.
  \end{itemize}
\end{resultados}

\section{De los algoritmos a los programas: estructura y organización del código}

Un algoritmo describe un procedimiento de solución con independencia de un lenguaje específico. Un programa expresa dicho procedimiento mediante instrucciones que pueden ser interpretadas o ejecutadas por un sistema computacional. La implementación exige respetar la sintaxis del lenguaje y conservar el significado de las operaciones definidas durante el diseño. \textcite[caps.~2 y 3]{farrell2013} relaciona la calidad de un programa con la claridad de sus instrucciones, la modularización y el empleo de estructuras de control.

En Python, la sangría delimita los bloques y \Codigo{\#} introduce un comentario \parencite{python312tutorial}. Organice los primeros programas en tres pasos: adquisición de datos, cálculo y presentación de resultados. \textcite{sorva2013} recomienda hacer explícito qué sucede al ejecutar el código. Para practicarlo, anote el valor de cada variable tras su asignación y anticipe qué mostrará la última línea.

\Ejemplo{Del modelo EPS al código}

Se desea calcular el área de un rectángulo de base 8 y altura 3, ambas en centímetros. El algoritmo recibe las dimensiones, multiplica sus valores y comunica el área. La versión siguiente usa datos fijos para comprobar primero la lógica.

\Python{ejemplos/03-01-area.py}

\textbf{Resultado esperado.} Área: 24 cm².

\begin{tablalibro}{2}{L{0.51}L{0.49}}
\caption{Correspondencia entre diseño e implementación}\label{tab:correspondencia-entre-diseno-e-implementacion}\\
\toprule
\rowcolor{Papel}
\CabeceraTabla{Diseño algorítmico} & \CabeceraTabla{Python} \\
\midrule
\endfirsthead
\SigueTabla{2}
\toprule
\rowcolor{Papel}
\CabeceraTabla{Diseño algorítmico} & \CabeceraTabla{Python} \\
\midrule
\endhead
\midrule
\endfoot
\bottomrule
\endlastfoot

Leer o establecer base y altura & \Codigo{base = 8; altura = 3} \\
Calcular área & \Codigo{area = base * altura} \\
Mostrar resultado & \Codigo{print(area)} \\
\end{tablalibro}


Cambie la base de 8 a 10, anticipe el área y compare el resultado al volver a ejecutar el ejemplo.

\section{Variables, constantes y tipos de datos}

Una variable es un nombre asociado a un valor durante la ejecución. En Python, la asignación vincula un nombre a un objeto sin exigir una declaración previa del tipo. Una constante representa un dato que debe conservar su valor durante el cálculo. PEP~8 recomienda escribir su nombre en mayúsculas como convención de lectura \parencite{pep8}. \textcite[cap.~2]{farrell2013} distingue los valores literales, las variables y las constantes nombradas como elementos del diseño de programas.

\begin{tablalibro}{3}{L{0.21}L{0.31}L{0.48}}
\caption{Tipos de datos básicos en Python}\label{tab:tipos-de-datos-basicos-en-python}\\
\toprule
\rowcolor{Papel}
\CabeceraTabla{Tipo} & \CabeceraTabla{Ejemplo} & \CabeceraTabla{Uso} \\
\midrule
\endfirsthead
\SigueTabla{3}
\toprule
\rowcolor{Papel}
\CabeceraTabla{Tipo} & \CabeceraTabla{Ejemplo} & \CabeceraTabla{Uso} \\
\midrule
\endhead
\midrule
\ContinuaTabla{3}
\endfoot
\bottomrule
\multicolumn{3}{@{}p{\linewidth}@{}}{\NotaLibro{Elaboración propia.}}\\
\endlastfoot

\Codigo{int} & \Codigo{12} & Cantidades enteras \\
\Codigo{float} & \Codigo{3.75} & Medidas aproximadas \\
\Codigo{str} & \Codigo{"Cuaderno"} & Texto y códigos \\
\Codigo{bool} & \Codigo{True} & Condiciones lógicas \\
\Codigo{NoneType} & \Codigo{None} & Ausencia de valor \\
\end{tablalibro}


\Ejemplo{Datos de un pedido}

Las cantidades son enteras; el nombre del producto es texto y el estado de disponibilidad es lógico. Para el dinero se emplean céntimos enteros, evitando errores de representación binaria de decimales.

\Python{ejemplos/03-02-datos-pedido.py}

\textbf{Resultado esperado.} Cuaderno 1350 \Codigo{True}.

Los identificadores no pueden comenzar por un dígito ni contener espacios. Se recomienda snake\_case, por ejemplo cantidad\_total. Los nombres distinguen mayúsculas y minúsculas. La función type() permite examinar el tipo de un valor, mientras \Codigo{int}(), \Codigo{float}() y \Codigo{str}() convierten valores compatibles. No todos los textos representan números: \Codigo{int}("tres") genera una excepción \Codigo{ValueError}.

Buena práctica. No confunda una representación textual con un valor numérico. El código de estudiante "00123" debe conservarse como texto si sus ceros iniciales son significativos.

\section{Operadores, expresiones y evaluación de resultados}

Una expresión combina valores, nombres y operadores para producir un resultado. Los operadores aritméticos permiten calcular; los relacionales comparan valores y producen resultados booleanos; los lógicos combinan condiciones. La asignación = no equivale a la comparación ==. Las operaciones deben corresponder a las reglas del problema, no únicamente a las posibilidades del lenguaje.

\begin{tablalibro}{3}{L{0.29}L{0.30}L{0.41}}
\caption{Operadores de uso frecuente}\label{tab:operadores-de-uso-frecuente}\\
\toprule
\rowcolor{Papel}
\CabeceraTabla{Categoría} & \CabeceraTabla{Operadores} & \CabeceraTabla{Ejemplo y resultado} \\
\midrule
\endfirsthead
\SigueTabla{3}
\toprule
\rowcolor{Papel}
\CabeceraTabla{Categoría} & \CabeceraTabla{Operadores} & \CabeceraTabla{Ejemplo y resultado} \\
\midrule
\endhead
\midrule
\endfoot
\bottomrule
\multicolumn{3}{@{}p{\linewidth}@{}}{\NotaLibro{Elaboración propia.}}\\
\endlastfoot

Aritmética & \Codigo{+  -  *  /} & \Codigo{7 / 2} \(\to\) \Codigo{3.5} \\
División entera y resto & \Codigo{//  \%} & \Codigo{7 // 2} \(\to\) \Codigo{3; 7 \% 2 \ensuremath{\rightarrow} 1} \\
Potencia & \Codigo{**} & \Codigo{2 ** 3} \(\to\) \Codigo{8} \\
Comparación & \Codigo{== != < <= > >=} & \Codigo{5 >= 3} \(\to\) \Codigo{True} \\
Lógica & \Codigo{and, or, not} & \Codigo{not True} \(\to\) \Codigo{False} \\
Asignación & \Codigo{=} & \Codigo{total = 10} \\
Actualización & \Codigo{+=, -=} & \Codigo{total += 5} \\
\end{tablalibro}


La precedencia determina qué operación se realiza primero. En expresiones usuales, la potencia precede a multiplicación y división, estas preceden a suma y resta, y las comparaciones se evalúan antes de las combinaciones booleanas. Dentro de los operadores lógicos, \Codigo{not} precede a \Codigo{and} y este a \Codigo{or}. Los paréntesis permiten explicitar la intención y reducir ambigüedades. Debe evitarse la división por cero.

\Ejemplo{Promedio ponderado}

Tres evaluaciones tienen pesos de 20 \%, 30 \% y 50 \%. Se calcula el promedio sin redondearlo antes de decidir la aprobación. El umbral 11 es únicamente una regla ficticia del ejercicio.

\Python{ejemplos/03-03-promedio.py}

\textbf{Resultado esperado.} Promedio: 14.40; Aprobado: \Codigo{True}.

En Python, \Codigo{and} y \Codigo{or} aplican evaluación de cortocircuito. Si una condición ya determina el resultado, la segunda puede no evaluarse. Por ejemplo, \Codigo{divisor != 0 and total / divisor > 5} evita calcular una división inválida cuando el divisor es cero. Esta propiedad no sustituye el análisis del dominio de los datos.

\section{Entrada y salida de información}

La entrada permite obtener datos del usuario o de otras fuentes. La salida comunica resultados, mensajes de error o estados del proceso. En una aplicación de consola, \Codigo{input()} recibe una línea de texto y \Codigo{print()} presenta información. Por tanto, la conversión de tipos es necesaria antes de realizar cálculos con una entrada numérica. Un mensaje de solicitud debe indicar qué dato se espera y en qué unidad.

\Ejemplo{Conversión de minutos}

Se solicita una cantidad entera de minutos no negativa y se convierte a horas y minutos. La captura de \Codigo{ValueError} distingue un formato inválido de un número fuera del dominio.

\Python{ejemplos/03-04-conversion-minutos.py}

\textbf{Resultado esperado.} Si se ingresa 135: 2 h y 15 min.

La instrucción \Codigo{try} permite intentar una operación que podría fallar y \Codigo{except} \Codigo{ValueError} gestiona este tipo particular de error. Se estudiará la depuración con mayor detalle en el \RefCapitulo{cap:04-pruebas-y-mejora}. La cadena formateada \Codigo{f"\{valor:.2f\}"} muestra dos cifras decimales, pero no altera el valor almacenado. Diseñar salidas claras es parte de la solución: un número sin etiqueta ni unidad puede ser interpretado incorrectamente.

\section{Estructuras secuenciales: ejecución ordenada de instrucciones}

Una secuencia ejecuta instrucciones una después de otra en el orden establecido. Se utiliza cuando todas las operaciones del proceso deben realizarse sin elegir entre alternativas ni repetir pasos. \textcite[cap.~3]{farrell2013} presenta la secuencia, la selección y la repetición como estructuras básicas de la programación estructurada. La dependencia entre operaciones determina el orden correcto de la secuencia.

\Ejemplo{Consumo eléctrico}

Se estima un importe ficticio a partir de dos lecturas de un medidor en kWh y un precio de 75 céntimos por kWh. Se supone que la lectura final es mayor o igual a la inicial.

\Python{ejemplos/03-05-consumo.py}

\textbf{Resultado esperado.} Consumo: 28 kWh; Importe: S/ 21.00.

\begin{tablalibro}{3}{R{0.16}L{0.50}L{0.34}}
\caption{Seguimiento de la secuencia}\label{tab:seguimiento-de-la-secuencia}\\
\toprule
\rowcolor{Papel}
\CabeceraTabla{Paso} & \CabeceraTabla{Variable modificada} & \CabeceraTabla{Valor} \\
\midrule
\endfirsthead
\SigueTabla{3}
\toprule
\rowcolor{Papel}
\CabeceraTabla{Paso} & \CabeceraTabla{Variable modificada} & \CabeceraTabla{Valor} \\
\midrule
\endhead
\midrule
\ContinuaTabla{3}
\endfoot
\bottomrule
\multicolumn{3}{@{}p{\linewidth}@{}}{\NotaLibro{Elaboración propia.}}\\
\endlastfoot

1 & lectura\_inicial & 120 \\
2 & lectura\_final & 148 \\
3 & TARIFA\_CENTIMOS & 75 \\
4 & consumo & 28 \\
5 & importe & 2100 \\
\end{tablalibro}


El cálculo de importe depende de consumo; invertir estas instrucciones podría utilizar un nombre todavía no definido. Una secuencia no puede resolver por sí sola la condición de una lectura final menor que la inicial: para ello se requiere una decisión.

\section{Estructuras selectivas: decisiones y condiciones lógicas}

Una estructura selectiva determina qué instrucciones ejecutar según el resultado de una condición. La forma simple \Codigo{if} ejecuta un bloque si la condición es verdadera. La forma \Codigo{if–else} dispone de dos alternativas, y la forma \Codigo{if–elif–else} evalúa alternativas ordenadas, ejecutando únicamente la primera condición verdadera o el bloque final. \textcite[cap.~4]{farrell2013} desarrolla la relación entre expresiones booleanas, operadores relacionales y selección.

\Ejemplo{Clasificación de existencias}

Una tienda considera agotado un artículo con cero unidades, bajo stock de una a cinco unidades y normal por encima de cinco. Las cantidades negativas son inválidas.

\Python{ejemplos/03-06-inventario.py}

\textbf{Resultado esperado.} Stock bajo.

\begin{tablalibro}{2}{R{0.32}L{0.68}}
\caption{Pruebas de límites de la selección}\label{tab:pruebas-de-limites-de-la-seleccion}\\
\toprule
\rowcolor{Papel}
\CabeceraTabla{Stock} & \CabeceraTabla{Resultado esperado} \\
\midrule
\endfirsthead
\SigueTabla{2}
\toprule
\rowcolor{Papel}
\CabeceraTabla{Stock} & \CabeceraTabla{Resultado esperado} \\
\midrule
\endhead
\midrule
\ContinuaTabla{2}
\endfoot
\bottomrule
\multicolumn{2}{@{}p{\linewidth}@{}}{\NotaLibro{Elaboración propia.}}\\
\endlastfoot

-1 & Dato inválido \\
0 & Agotado \\
1 & Stock bajo \\
5 & Stock bajo \\
6 & Stock normal \\
\end{tablalibro}


El orden de las condiciones importa: si se comprobara stock <= 5 antes de descartar los valores negativos y el cero, estos podrían clasificarse erróneamente como stock bajo. Cuando varias reglas deben aplicarse independientemente, se utilizan instrucciones \Codigo{if} separadas en lugar de \Codigo{elif}. Las condiciones compuestas permiten expresar rangos con \Codigo{and}, alternativas con \Codigo{or} y negaciones con \Codigo{not}.

Error frecuente. Escribir \Codigo{if} nota >= 0 \Codigo{and} <= 20 no es válido. Use 0 <= nota <= 20 o nota >= 0 \Codigo{and} nota <= 20.

\section{Estructuras repetitivas: ciclos, contadores y acumuladores}

Un ciclo repite un bloque de instrucciones. El ciclo \Codigo{for} de Python recorre los elementos de un iterable, como una lista o un rango; \Codigo{while} repite mientras una condición sea verdadera. Un contador registra cuántas veces ocurre un evento, mientras un acumulador integra valores sucesivos, por ejemplo una suma. \textcite[cap.~5]{farrell2013} analiza el control de ciclos y la acumulación de totales, así como errores de inicialización y actualización.

\Ejemplo{Registro de gastos}

Se dispone de cuatro gastos en céntimos. Se calcula su total y se cuenta cuántos superan diez soles.

\Python{ejemplos/03-07-registro-gastos.py}

\textbf{Resultado esperado.} Total: S/ 40.00; Gastos mayores a S/ 10: 2.

\begin{tablalibro}{4}{R{0.25}R{0.25}R{0.25}R{0.25}}
\caption{Prueba de escritorio del ciclo}\label{tab:prueba-de-escritorio-del-ciclo}\\
\toprule
\rowcolor{Papel}
\CabeceraTabla{Iteración} & \CabeceraTabla{Gasto} & \CabeceraTabla{Total} & \CabeceraTabla{Mayores} \\
\midrule
\endfirsthead
\SigueTabla{4}
\toprule
\rowcolor{Papel}
\CabeceraTabla{Iteración} & \CabeceraTabla{Gasto} & \CabeceraTabla{Total} & \CabeceraTabla{Mayores} \\
\midrule
\endhead
\midrule
\ContinuaTabla{4}
\endfoot
\bottomrule
\multicolumn{4}{@{}p{\linewidth}@{}}{\NotaLibro{Elaboración propia.}}\\
\endlastfoot

Inicial & — & 0 & 0 \\
1 & 1200 & 1200 & 1 \\
2 & 800 & 2000 & 1 \\
3 & 1500 & 3500 & 2 \\
4 & 500 & 4000 & 2 \\
\end{tablalibro}


\Ejemplo{Repetición condicionada}

Se muestra una cuenta regresiva. El valor de control disminuye en cada iteración y alcanza la condición de salida.

\Python{ejemplos/03-08-repeticion.py}

\textbf{Resultado esperado.} 3, 2, 1 y Fin (en líneas independientes).

\Codigo{range}(1, 4) genera 1, 2 y 3: su límite superior es excluyente. En \Codigo{while} debe comprobarse que exista una condición de salida alcanzable. Un ciclo puede ejecutar cero veces si la condición inicial es falsa. La instrucción \Codigo{break} permite finalizar anticipadamente y continue pasa a la siguiente iteración; conviene utilizarlas cuando mantienen clara la lógica del proceso.

\section{Modularización de soluciones mediante funciones y procedimientos}

Modularizar significa dividir la implementación en unidades con responsabilidades definidas. Una función recibe argumentos, realiza una tarea y puede devolver un resultado mediante \Codigo{return}. En Python no existe una construcción separada denominada procedimiento: una función que se utiliza principalmente por sus efectos y no devuelve explícitamente un valor retorna \Codigo{None}. \textcite[cap.~2]{farrell2013} destaca la abstracción y la reutilización como ventajas de la modularización.

\Ejemplo{Funciones para calcular un pedido}

El cálculo se separa de la presentación. La función subtotal puede utilizarse en diferentes pedidos sin repetir la operación.

\Python[lastline=8]{ejemplos/03-09-funciones.py}

\textbf{Resultado esperado.} Importe: S/ 14.00.

Los parámetros son los nombres definidos en la firma de una función; los argumentos son los valores proporcionados al invocarla. Las variables creadas dentro de una función son normalmente locales. Return permite que el resultado se utilice en otros cálculos; \Codigo{print} solo lo presenta. La interfaz de una función debe indicar qué datos admite y qué resultado devuelve. Evitar dependencias innecesarias de variables globales facilita las pruebas.

\Python[firstline=11]{ejemplos/03-09-funciones.py}

Estas comprobaciones expresan resultados esperados para la función anterior. Si una igualdad es falsa, \Codigo{assert} produce \Codigo{AssertionError} en la ejecución normal. No sustituye la validación de datos de usuario; sirve aquí como introducción a pruebas de unidades pequeñas.

\section{Organización y manipulación básica de colecciones de datos}

Una colección agrupa varios valores. Permite representar conjuntos de registros, recorrerlos y efectuar cálculos sin crear una variable independiente para cada dato. La elección de la colección depende de las operaciones necesarias: conservar orden, acceder por posición, asociar claves o evitar duplicados. En Python existen listas, tuplas, diccionarios y conjuntos; sus propiedades deben distinguirse antes de elegir una estructura.

\begin{tablalibro}{3}{L{0.22}L{0.48}L{0.30}}
\caption{Colecciones básicas}\label{tab:colecciones-basicas}\\
\toprule
\rowcolor{Papel}
\CabeceraTabla{Colección} & \CabeceraTabla{Propiedades} & \CabeceraTabla{Ejemplo} \\
\midrule
\endfirsthead
\SigueTabla{3}
\toprule
\rowcolor{Papel}
\CabeceraTabla{Colección} & \CabeceraTabla{Propiedades} & \CabeceraTabla{Ejemplo} \\
\midrule
\endhead
\midrule
\endfoot
\bottomrule
\endlastfoot

Lista & Ordenada, modificable, admite repetidos & \Codigo{[12, 8, 12]} \\
Tupla & Ordenada, estructura no modificable & \Codigo{(12, 8)} \\
Diccionario & Asocia claves únicas a valores & \Codigo{\{"A01": 12\}} \\
Conjunto & Elementos únicos, sin orden garantizado & \Codigo{\{12, 8\}} \\
\end{tablalibro}


\Ejemplo{Lista de calificaciones}

Se incorpora una calificación y se calcula el promedio si la lista contiene elementos. El índice inicial de una lista es cero.

\Python{ejemplos/03-10-lista-notas.py}

\textbf{Resultado esperado.} Promedio: 13.00; Primera nota: 13.

\Ejemplo{Diccionario de inventario}

Se actualiza la cantidad de un artículo por su código y se recorren las parejas de clave y valor. La función get permite establecer un valor por defecto si una clave no existe.

\Python{ejemplos/03-11-diccionario.py}

\textbf{Resultado esperado.} A01 10; A02 3; A03 5; 0 (en líneas independientes).

Modificar una colección mientras se recorre puede provocar errores o resultados inesperados. Cuando se necesita filtrar elementos, suele resultar más claro construir una nueva lista. Las colecciones vacías requieren especial atención: no es posible calcular un promedio sin elementos ni acceder a una posición inexistente. Estas situaciones se incorporan al diseño de pruebas.

\SeccionLibro[Caso práctico]{Caso práctico. Implementación progresiva de una solución}

Se implementa el algoritmo del \RefCapitulo{cap:02-diseno-algoritmico} con los mismos materiales y un presupuesto de S/ 100,00. En Python, la lista de registros determina la cantidad de elementos que se recorrerán. El saldo conserva su signo para indicar el importe disponible o el faltante.

\begin{tablalibro}{4}{L{0.35}R{0.14}R{0.26}R{0.25}}
\caption{Datos del escenario de trabajo}\label{tab:datos-del-escenario-de-trabajo}\\
\toprule
\rowcolor{Papel}
\CabeceraTabla{Material} & \CabeceraTabla{Cantidad} & \CabeceraTabla{Precio unitario (céntimos)} & \CabeceraTabla{Subtotal (céntimos)} \\
\midrule
\endfirsthead
\SigueTabla{4}
\toprule
\rowcolor{Papel}
\CabeceraTabla{Material} & \CabeceraTabla{Cantidad} & \CabeceraTabla{Precio unitario (céntimos)} & \CabeceraTabla{Subtotal (céntimos)} \\
\midrule
\endhead
\midrule
\endfoot
\bottomrule
\multicolumn{4}{@{}p{\linewidth}@{}}{\NotaLibro{Elaboración propia.}}\\
\endlastfoot

Cuadernos & 5 & 450 & 2250 \\
Lápices & 10 & 120 & 1200 \\
Cartulinas & 4 & 300 & 1200 \\
\textbf{Total} & \textbf{—} & \textbf{—} & \textbf{4650} \\
\end{tablalibro}


\subsection{Etapa 1. Variables y operaciones}

Con un solo producto se comprueba la fórmula del subtotal y la representación monetaria. Los listados de las etapas 1 a 3 se ejecutan en orden.

\Python[firstline=1,lastline=5]{ejemplos/03-caso-etapas.py}

\subsection{Etapa 2. Decisiones}

Se clasifica el importe mediante una condición. La igualdad corresponde a una compra que consume todo el presupuesto, sin exceso.

\Python[firstline=7,lastline=10]{ejemplos/03-caso-etapas.py}

\subsection{Etapa 3. Ciclos y colecciones}

Se representan los materiales como registros y se acumulan sus subtotales. Esta versión supone datos válidos.

\Python[firstline=12,lastline=18]{ejemplos/03-caso-etapas.py}

\subsection{Etapa 4. Funciones y programa integrado}

Cada elemento de la lista es una terna con nombre, cantidad y precio en céntimos. \Codigo{evaluar\_compra(presupuesto, materiales)} devuelve \Codigo{(total, saldo, viable)}. Los dos primeros valores se expresan en céntimos y el tercero es booleano. Si un dato incumple las reglas de la \RefTabla{tab:requerimientos-del-caso-practico}, la función lanza \Codigo{ValueError}. El programa separa este cálculo de la impresión del informe.

\Python*{ejemplos/03-caso-integrado.py}

La condición final ejecuta \Codigo{main} al iniciar directamente el archivo. Así es posible importar la función para probarla sin imprimir el informe. \Codigo{divmod} separa soles y céntimos mediante el cociente y el resto de la división entre 100.

\subsection{Etapa 5. Ejecución y comprobación}

\Salida{03-caso-integrado.txt}

\begin{tablalibro}{3}{R{0.24}L{0.36}L{0.40}}
\caption{Casos de prueba del programa}\label{tab:casos-de-prueba-del-programa}\\
\toprule
\rowcolor{Papel}
\CabeceraTabla{Presupuesto} & \CabeceraTabla{Materiales (cantidad × precio)} & \CabeceraTabla{Resultado esperado} \\
\midrule
\endfirsthead
\SigueTabla{3}
\toprule
\rowcolor{Papel}
\CabeceraTabla{Presupuesto} & \CabeceraTabla{Materiales (cantidad × precio)} & \CabeceraTabla{Resultado esperado} \\
\midrule
\endhead
\midrule
\endfoot
\bottomrule
\multicolumn{3}{@{}p{\linewidth}@{}}{\NotaLibro{Elaboración propia.}}\\
\endlastfoot

10000 & 5×450; 10×120; 4×300 & Viable; saldo 5350 \\
4650 & Mismos materiales & Viable; saldo 0 \\
4000 & Mismos materiales & Insuficiente; exceso 650 \\
10000 & Lista vacía & Viable; saldo 10000 \\
10000 & 0×450 & Error: cantidad inválida \\
-1 & Lista vacía & Error: presupuesto inválido \\
\end{tablalibro}


Compare el total y el saldo con la prueba de escritorio del capítulo anterior. Después, cambie el presupuesto a 4650 y a 4000 céntimos y anticipe el saldo y la viabilidad. El \RefCapitulo{cap:04-pruebas-y-mejora} automatiza estas comprobaciones y vincula los casos de prueba con sus requisitos.

\SeccionLibro{Cierre del capítulo}

\subsection{Resumen}

La programación convierte algoritmos en instrucciones ejecutables. Las variables, constantes y tipos representan los datos; los operadores expresan cálculos y condiciones; la entrada y salida conectan la solución con el usuario. Las estructuras secuenciales, selectivas y repetitivas organizan la ejecución. Las funciones permiten dividir responsabilidades, reutilizar procedimientos y comprobar resultados de forma aislada, mientras las colecciones agrupan datos y hacen posible procesar cantidades variables de información. El caso integrador mostró cómo construir una solución progresivamente y contrastarla con su diseño inicial.

\Needspace{11\baselineskip}
\subsection{Términos clave}

\begin{tablalibro}{2}{L{0.31}L{0.69}}
\caption{Glosario del capítulo}\label{tab:glosario-del-capitulo-3}\\
\toprule
\rowcolor{Papel}
\CabeceraTabla{Término} & \CabeceraTabla{Definición} \\
\midrule
\endfirsthead
\SigueTabla{2}
\toprule
\rowcolor{Papel}
\CabeceraTabla{Término} & \CabeceraTabla{Definición} \\
\midrule
\endhead
\midrule
\endfoot
\bottomrule
\multicolumn{2}{@{}p{\linewidth}@{}}{\NotaLibro{Elaboración propia.}}\\
\endlastfoot

Programa & Implementación ejecutable de un procedimiento. \\
Variable & Nombre asociado a un valor. \\
Constante & Valor que se pretende mantener sin cambios. \\
Tipo de dato & Categoría que determina operaciones compatibles. \\
Expresión & Combinación que produce un valor. \\
Selección & Ejecución de una alternativa según una condición. \\
Ciclo & Repetición de instrucciones. \\
Contador & Variable que registra ocurrencias. \\
Acumulador & Variable que integra valores sucesivos. \\
Función & Unidad invocable que puede devolver un resultado. \\
Parámetro & Nombre que recibe un argumento. \\
Colección & Estructura que agrupa valores. \\
Sangría & Espacio inicial que delimita bloques en Python. \\
Prueba de escritorio & Seguimiento manual de estados y resultados. \\
\end{tablalibro}


\subsection{Preguntas de comprensión}

\begin{enumerate}
\item ¿Qué decisiones quedan explícitas al llevar un algoritmo a Python?
\item ¿Por qué Python distingue un entero de una cadena de texto?
\item ¿Qué diferencia existe entre = y ==?
\item ¿Cómo influye la precedencia en una expresión?
\item ¿Por qué \Codigo{input()} requiere conversiones para realizar cálculos?
\item ¿Cuándo es suficiente una estructura secuencial?
\item ¿Qué diferencia existe entre \Codigo{if} independientes e \Codigo{if–elif–else}?
\item ¿Cómo se garantiza que un ciclo \Codigo{while} termine?
\item ¿En qué se diferencian un contador y un acumulador?
\item ¿Qué distingue \Codigo{return} de \Codigo{print}?
\item ¿Cuándo conviene elegir un diccionario en lugar de una lista?
\item ¿Por qué deben probarse colecciones vacías y valores límite?
\end{enumerate}

\subsection{Ejercicios de construcción de programas}

\begin{enumerate}
\item Conversión de unidades. Lea segundos enteros no negativos y muestre horas, minutos y segundos. Pruebe 0, 59, 60 y 3601.
\item Calificaciones. Solicite tres notas entre 0 y 20, calcule su promedio y determine si alcanza un umbral ficticio de 11. Rechace entradas inválidas.
\item Inventario. Lea existencias y unidades solicitadas; compruebe disponibilidad y calcule las existencias restantes. Pruebe la igualdad y cantidades negativas.
\item Registro de consumo. Recorra siete consumos diarios, calcule total, promedio y máximo. Defina qué ocurre ante datos inválidos.
\item Contador de asistencia. Recorra una lista de marcas P y A; calcule presentes, ausentes y porcentaje de asistencia. Trate la lista vacía.
\item Funciones. Defina una función que calcule el importe de un pedido en céntimos y otra que lo formatee en soles. Compruebe al menos tres entradas.
\item Colecciones. Cree un diccionario de productos y cantidades. Permita actualizar existencias y consultar un código inexistente sin provocar un error.
\item Filtrado. Construya una nueva lista con las calificaciones aprobatorias de una lista original, sin modificar esta última.
\end{enumerate}

\subsection{Desafíos de programación progresiva}

\textbf{Desafío 1.} Biblioteca escolar. Implemente primero el registro de un préstamo, luego una decisión de vencimiento basada en días transcurridos, después un ciclo para varios préstamos y finalmente funciones para calcular y presentar un informe. Incluya pruebas de cero días y del límite de préstamo.

\textbf{Desafío 2.} Monitoreo ambiental. Registre mediciones de temperatura válidas, calcule promedio y máximo, cuente las que superan un umbral configurable y encapsule el procesamiento en funciones. Establezca una política para la lista vacía. Compare los resultados con una tabla manual.

\textbf{Desafío 3.} Control de presupuesto ampliado. Extienda el caso integrador para leer materiales desde la consola, detectar códigos duplicados y mostrar el producto de mayor subtotal. Defina cómo resolver los empates y terminar la entrada de datos. Compruebe ambas reglas con casos de prueba.
